\documentclass[12pt,a4paper]{article}

% Paquetes necesarios
\usepackage[utf8]{inputenc}
\usepackage[spanish]{babel}
\usepackage[margin=2.5cm]{geometry}
\usepackage{amsmath}
\usepackage{amsfonts}
\usepackage{amssymb}
\usepackage{graphicx}
\usepackage{xcolor}
\usepackage{fancyhdr}
\usepackage{titlesec}
\usepackage{tocloft}
\usepackage{enumitem}
\usepackage{multicol}
\usepackage{mdframed}
\usepackage{setspace}
\usepackage{float}
\usepackage{hyperref}

% Configuración de hyperref
\hypersetup{
    colorlinks=true,
    linkcolor=forestgreen,
    filecolor=magenta,      
    urlcolor=cyan,
    citecolor=forestgreen,
    pdfborderstyle={/S/U/W 1}
}

% Definir colores del tema
\definecolor{forestgreen}{RGB}{27,67,50}
\definecolor{sagegreen}{RGB}{45,90,39}
\definecolor{olivegreen}{RGB}{64,145,108}
\definecolor{mintgreen}{RGB}{82,183,136}
\definecolor{lightmint}{RGB}{116,198,157}
\definecolor{creammint}{RGB}{183,228,199}
\definecolor{softcream}{RGB}{216,243,220}
\definecolor{paperwhite}{RGB}{248,253,248}
\definecolor{ambergold}{RGB}{214,158,46}

% Configuración de títulos de sección
\titleformat{\section}
{\color{forestgreen}\Large\bfseries}
{\color{forestgreen}\thesection.}{1em}{}

\titleformat{\subsection}
{\color{olivegreen}\large\bfseries}
{\color{olivegreen}\thesubsection.}{1em}{}

\titleformat{\subsubsection}
{\color{sagegreen}\normalsize\bfseries}
{\color{sagegreen}\thesubsubsection.}{1em}{}

% Configuración de encabezados y pies de página
\pagestyle{fancy}
\fancyhf{}
\fancyhead[L]{\textcolor{forestgreen}{\textbf{Análisis Exegético Profundo}}}
\fancyhead[R]{\textcolor{olivegreen}{[Tu Nombre]}}
\fancyfoot[C]{\textcolor{sagegreen}{\thepage}}
\renewcommand{\headrulewidth}{1pt}
\renewcommand{\footrulewidth}{0.5pt}
\renewcommand{\headrule}{\hbox to\headwidth{\color{forestgreen}\leaders\hrule height \headrulewidth\hfill}}
\renewcommand{\footrule}{\hbox to\headwidth{\color{olivegreen}\leaders\hrule height \footrulewidth\hfill}}

% Configuración del espaciado
\onehalfspacing
\setlength{\parindent}{0pt}
\setlength{\parskip}{0.5em}

% Comandos personalizados
\newcommand{\versiculoclave}[1]{\textcolor{olivegreen}{\textbf{#1}}}
\newcommand{\palabraclave}[1]{\textcolor{ambergold}{\textbf{#1}}}
\newcommand{\principio}[1]{\textcolor{forestgreen}{\textbf{#1}}}
\newcommand{\aplicacion}[1]{\textcolor{sagegreen}{\textit{#1}}}

% Entornos personalizados para cajas de información
\mdfdefinestyle{InfoBox}{
    backgroundcolor=softcream,
    linecolor=olivegreen,
    linewidth=2pt,
    roundcorner=5pt,
    innertopmargin=10pt,
    innerbottommargin=10pt,
    innerrightmargin=15pt,
    innerleftmargin=15pt,
    skipabove=10pt,
    skipbelow=10pt
}

\mdfdefinestyle{AlertBox}{
    backgroundcolor=creammint,
    linecolor=ambergold,
    linewidth=2pt,
    roundcorner=5pt,
    innertopmargin=10pt,
    innerbottommargin=10pt,
    innerrightmargin=15pt,
    innerleftmargin=15pt,
    skipabove=10pt,
    skipbelow=10pt
}

\mdfdefinestyle{ImportantBox}{
    backgroundcolor=lightmint,
    linecolor=forestgreen,
    linewidth=3pt,
    roundcorner=5pt,
    innertopmargin=15pt,
    innerbottommargin=15pt,
    innerrightmargin=20pt,
    innerleftmargin=20pt,
    skipabove=15pt,
    skipbelow=15pt
}

% Información del documento
\title{\color{forestgreen}\Huge\textbf{Análisis Exegético Profundo}\\
\vspace{0.5cm}
\color{olivegreen}\Large [Escribir aquí el título específico de tu texto]}
\author{\color{sagegreen}\large [Tu Nombre Completo]\\
\textcolor{forestgreen}{Escuela de Ministerio Doctrinal}\\
\textcolor{olivegreen}{Ministerio ``La Gloria es del Señor''}}
\date{\color{ambergold}\large\today}

\begin{document}

% Portada
\maketitle
\thispagestyle{empty}

\vspace{2cm}
\begin{center}
\begin{mdframed}[style=ImportantBox]
\centering
\textcolor{forestgreen}{\Large\textbf{Proyecto de Investigación Bíblica}}\\
\vspace{0.5cm}
\textcolor{olivegreen}{\large Doctrina Intermedia}\\
\vspace{0.5cm}
\textcolor{sagegreen}{Un estudio sistemático y riguroso de la Palabra de Dios}
\end{mdframed}
\end{center}

\newpage

% Tabla de contenidos
\tableofcontents
\newpage

% ====================
% SECCIÓN I: INTRODUCCIÓN
% ====================
\section{Introducción}

\subsection{El Texto Seleccionado}

\begin{mdframed}[style=InfoBox]
\textbf{Pasaje Bíblico:} \versiculoclave{[Cita bíblica completa]}

\vspace{0.5cm}
[Escribir aquí el texto completo del pasaje en una traducción fiel y clara]
\end{mdframed}

\begin{mdframed}[style=AlertBox]
\textbf{¿Por qué este texto?}
\begin{itemize}[leftmargin=1.5em, rightmargin=0.5em]
    \item Razón 1 para la selección del texto
    \item Razón 2 para la selección del texto  
    \item Relevancia específica para la congregación actual
    \item Importancia teológica o doctrinal del pasaje
\end{itemize}
\end{mdframed}

\subsection{Importancia del Análisis Exegético}

La exégesis bíblica es el proceso sistemático de interpretar las Escrituras mediante el cual extraemos el significado original que el autor inspirado quiso comunicar a su audiencia. Este método riguroso nos permite:

\begin{itemize}[leftmargin=*]
    \item \textbf{Extraer el significado original:} Descubrir qué quería decir el texto en su contexto histórico original
    \item \textbf{Entender la intención divina:} Comprender el propósito de Dios al inspirar este pasaje específico
    \item \textbf{Aplicar metodología rigurosa:} Usar herramientas académicas y espirituales apropiadas
    \item \textbf{Evitar interpretaciones subjetivas:} Prevenir lecturas erróneas basadas en prejuicios personales
\end{itemize}

\begin{mdframed}[style=ImportantBox]
\textbf{Meta de este Estudio:} Descubrir juntos lo que Dios quería comunicar a través de este pasaje y cómo se aplica de manera práctica y transformadora a nuestras vidas contemporáneas.
\end{mdframed}

\subsection{Contexto Histórico del Pasaje}

\subsubsection{Trasfondo Histórico}

\begin{itemize}[leftmargin=*]
    \item \textbf{Período Histórico:} [Especificar la época histórica en que fue escrito]
    \item \textbf{Autor Humano:} [Información detallada sobre el autor y sus credenciales]
    \item \textbf{Audiencia Original:} [Descripción de los destinatarios inmediatos]
    \item \textbf{Circunstancias de Composición:} [Situación histórica que motivó la escritura del texto]
    \item \textbf{Propósito Inmediato:} [Razón específica por la cual se escribió este pasaje]
\end{itemize}

\subsubsection{Contexto Geográfico y Cultural}

\begin{multicols}{2}
\begin{itemize}[leftmargin=*]
    \item \textbf{Ubicación:} [Lugar geográfico específico]
    \item \textbf{Costumbres:} [Prácticas culturales relevantes]
    \item \textbf{Religión:} [Contexto religioso de la época]
    \item \textbf{Política:} [Situación política imperante]
    \item \textbf{Economía:} [Condiciones económicas]
    \item \textbf{Sociedad:} [Estructura social y dinámicas]
\end{itemize}
\end{multicols}

% ====================
% SECCIÓN II: ANÁLISIS DEL TEXTO
% ====================
\section{Análisis del Texto}

\subsection{Estructura Literaria}

\subsubsection{Género Literario}
\textbf{Clasificación:} [Narrativa/Poesía/Epístola/Profecía/Apocalíptica/Sabiduría/Legal/Otros]

El reconocimiento del género literario es fundamental para la interpretación correcta, ya que cada género tiene sus propias reglas hermenéuticas y características distintivas.

\subsubsection{Estructura del Pasaje}

\begin{enumerate}[leftmargin=*]
    \item \textbf{Versículos X-Y:} [Descripción detallada de la primera sección y su función literaria]
    \item \textbf{Versículos X-Y:} [Descripción detallada de la segunda sección y su desarrollo temático]  
    \item \textbf{Versículos X-Y:} [Descripción detallada de la tercera sección y su conclusión]
\end{enumerate}

\begin{mdframed}[style=AlertBox]
\textbf{Elementos Retóricos Identificados:}
[Mencionar y explicar paralelismos, quiasmos, inclusiones, repeticiones, aliteraciones, figuras retóricas, etc., y su propósito en el texto]
\end{mdframed}

\subsection{Palabras Clave en el Idioma Original}

\subsubsection{Análisis Léxico Detallado}

\textbf{1. \palabraclave{[Primera Palabra Clave]}}
\begin{itemize}[leftmargin=*]
    \item \textbf{Idioma Original:} [Griego/Hebreo/Arameo]
    \item \textbf{Transliteración:} [Escritura fonética]
    \item \textbf{Número Strong:} [Código de concordancia]
    \item \textbf{Definición Básica:} [Significado fundamental de la palabra]
    \item \textbf{Matices Semánticos:} [Diferentes acepciones según el contexto]
    \item \textbf{Uso en Otros Pasajes:} [Referencias donde aparece la misma palabra]
\end{itemize}

\textbf{2. \palabraclave{[Segunda Palabra Clave]}}
\begin{itemize}[leftmargin=*]
    \item \textbf{Idioma Original:} [Griego/Hebreo/Arameo]
    \item \textbf{Transliteración:} [Escritura fonética]
    \item \textbf{Número Strong:} [Código de concordancia]
    \item \textbf{Definición Básica:} [Significado fundamental de la palabra]
    \item \textbf{Matices Semánticos:} [Diferentes acepciones según el contexto]
    \item \textbf{Uso en Otros Pasajes:} [Referencias donde aparece la misma palabra]
\end{itemize}

\textbf{3. \palabraclave{[Tercera Palabra Clave]}}
\begin{itemize}[leftmargin=*]
    \item \textbf{Idioma Original:} [Griego/Hebreo/Arameo]
    \item \textbf{Transliteración:} [Escritura fonética]
    \item \textbf{Número Strong:} [Código de concordancia]
    \item \textbf{Definición Básica:} [Significado fundamental de la palabra]
    \item \textbf{Matices Semánticos:} [Diferentes acepciones según el contexto]
    \item \textbf{Uso en Otros Pasajes:} [Referencias donde aparece la misma palabra]
\end{itemize}

\subsection{Análisis Gramatical}

\subsubsection{Elementos Gramaticales Fundamentales}

\begin{itemize}[leftmargin=*]
    \item \textbf{Tiempos Verbales:} [Análisis detallado de los verbos principales y su significado temporal/aspectual]
    \item \textbf{Conectores Lógicos:} [Examen de conjunciones, preposiciones y partículas que revelan las relaciones lógicas]
    \item \textbf{Estructura Sintáctica:} [Análisis de las relaciones entre cláusulas principales y subordinadas]
    \item \textbf{Modo y Voz Verbal:} [Significado de indicativo, subjuntivo, imperativo; activa, pasiva, media]
\end{itemize}

\begin{mdframed}[style=AlertBox]
\textbf{Flujo del Argumento:}
[Explicación detallada de cómo se desarrolla lógicamente el argumento del pasaje, mostrando la progresión del pensamiento del autor y las conexiones entre las ideas principales]
\end{mdframed}

\subsection{Contexto Canónico}

\subsubsection{Contexto Inmediato}

\begin{itemize}[leftmargin=*]
    \item \textbf{Contexto Anterior:} [Análisis de los versículos/capítulos precedentes y su relación con nuestro texto]
    \item \textbf{Contexto Posterior:} [Examen de lo que sigue al pasaje y cómo se conecta]
    \item \textbf{Flujo Temático del Libro:} [Cómo encaja este texto en el propósito general del libro bíblico]
\end{itemize}

\subsubsection{Textos Paralelos y Referencias Cruzadas}

\begin{itemize}[leftmargin=*]
    \item \textbf{[Referencia Bíblica 1]:} [Descripción detallada de la conexión temática, verbal o conceptual]
    \item \textbf{[Referencia Bíblica 2]:} [Explicación de cómo este texto ilumina o complementa nuestro pasaje]
    \item \textbf{[Referencia Bíblica 3]:} [Análisis de paralelos o contrastes significativos]
    \item \textbf{[Referencia Bíblica 4]:} [Desarrollo progresivo del tema en la revelación bíblica]
\end{itemize}

% ====================
% SECCIÓN III: MENSAJE CENTRAL
% ====================
\section{Mensaje Central}

\subsection{La Verdad Principal del Texto}

\subsubsection{Propósito del Autor Original}
[Desarrollar en 2-3 párrafos una explicación clara y precisa de lo que el autor inspirado quería comunicar a su audiencia inmediata, basándose en el análisis exegético realizado]

\begin{mdframed}[style=ImportantBox]
\textbf{Mensaje Central del Pasaje:}

\centering
\Large \principio{[Declaración clara, concisa y precisa del mensaje principal del texto, formulada en una oración completa que capture la esencia de la enseñanza bíblica]}
\end{mdframed}

\subsubsection{Puntos Teológicos Fundamentales}

\begin{enumerate}[leftmargin=*]
    \item \textbf{[Primer Punto Teológico]:} [Desarrollo detallado del primer aspecto doctrinal que emerge del texto]
    \item \textbf{[Segundo Punto Teológico]:} [Explicación del segundo elemento teológico significativo]
    \item \textbf{[Tercer Punto Teológico]:} [Elaboración del tercer componente doctrinal importante]
\end{enumerate}

\subsection{Contribución al Canon Bíblico}

\subsubsection{Desarrollo Progresivo de la Revelación}

\begin{itemize}[leftmargin=*]
    \item \textbf{Fundamento en el Antiguo Testamento:} [Cómo este texto se relaciona con y desarrolla temas establecidos en el AT]
    \item \textbf{Desarrollo en el Nuevo Testamento:} [Cómo se amplía, cumple o clarifica este tema en el NT]
    \item \textbf{Contribución Única:} [Qué aspecto distintivo o revelación especial aporta este texto específico]
    \item \textbf{Implicaciones Escatológicas:} [Cómo se relaciona con la consumación futura de los propósitos de Dios]
\end{itemize}

\begin{mdframed}[style=AlertBox]
\textbf{Relevancia Teológica:}
[Explicación comprensiva de por qué este texto es importante para la doctrina cristiana, qué vacío llena en nuestra comprensión de Dios, y cómo contribuye al sistema teológico general de las Escrituras]
\end{mdframed}

% ====================
% SECCIÓN IV: APLICACIÓN PRÁCTICA
% ====================
\section{Aplicación Práctica}

\subsection{Principios Universales}

\subsubsection{Principios Extraídos del Texto}

\begin{enumerate}[leftmargin=*]
    \item \principio{[Primer Principio Universal]:} [Explicación detallada de cómo este principio se deriva directamente del texto y trasciende las barreras culturales]
    \item \principio{[Segundo Principio Universal]:} [Desarrollo del segundo principio con ejemplos de su aplicabilidad transcultural]
    \item \principio{[Tercer Principio Universal]:} [Elaboración del tercer principio mostrando su relevancia contemporánea]
\end{enumerate}

\begin{mdframed}[style=InfoBox]
\textbf{Validación de los Principios:}

¿Cómo sabemos que estos principios son válidos para nosotros hoy?
\begin{itemize}[leftmargin=*]
    \item Están claramente basados en el significado original del texto bíblico
    \item Son consistentes con el resto de la revelación escritural
    \item Trascienden las particularidades culturales y temporales del texto original
    \item Abordan aspectos fundamentales de la naturaleza humana y divina
\end{itemize}
\end{mdframed}

\subsection{Aplicaciones Específicas}

\subsubsection{Para Toda la Congregación}
\aplicacion{[Desarrollar una aplicación general y práctica que sea relevante para todos los miembros de la congregación, con ejemplos concretos de cómo vivir este principio]}

\subsubsection{Para Líderes Espirituales}
\aplicacion{[Explicar cómo este texto debe informar y transformar el liderazgo cristiano, con aplicaciones específicas para pastores, ancianos y maestros]}

\subsubsection{Para las Familias}
\aplicacion{[Mostrar cómo estos principios se aplican en el contexto familiar, incluyendo relaciones matrimoniales, crianza de hijos y dinámicas del hogar cristiano]}

\subsubsection{Para la Vida Personal}
\aplicacion{[Desarrollar aplicaciones para el crecimiento espiritual individual, devocionales personal y transformación del carácter]}

\begin{mdframed}[style=AlertBox]
\textbf{Desafío Práctico Semanal:}
\aplicacion{[Formular un desafío concreto, específico y medible que los oyentes puedan implementar durante la semana siguiente]}
\end{mdframed}

\subsection{Ilustración Contemporánea}

[Narrar una ilustración contemporánea extensa y bien desarrollada que ayude a entender y aplicar el texto. Esta puede ser una historia real, una analogía detallada, un ejemplo de la cultura actual, una situación cotidiana, o una metáfora que ilumine el significado del pasaje]

\begin{mdframed}[style=InfoBox]
\textbf{Conexión con el Texto Bíblico:}
[Explicar claramente y en detalle cómo la ilustración se conecta con el mensaje del pasaje bíblico, qué aspectos ilumina, y por qué es efectiva para comunicar la verdad bíblica]
\end{mdframed}

% ====================
% SECCIÓN V: CONCLUSIÓN
% ====================
\section{Conclusión}

\subsection{Resumen de Hallazgos}

A través de este análisis exegético riguroso, hemos descubierto varias verdades fundamentales:

\begin{itemize}[leftmargin=*]
    \item \textbf{[Hallazgo Principal 1]:} [Resumir el primer descubrimiento significativo del estudio]
    \item \textbf{[Hallazgo Principal 2]:} [Sintetizar el segundo insight importante]
    \item \textbf{[Hallazgo Principal 3]:} [Recapitular el tercer aspecto crucial revelado]
\end{itemize}

\begin{mdframed}[style=ImportantBox]
\textbf{El Mensaje para Nosotros Hoy:}

\centering
\Large \principio{[Reafirmar el mensaje central del texto, pero ahora adaptado y aplicado específicamente para la audiencia contemporánea]}
\end{mdframed}

\subsection{Llamado a la Acción}

\subsubsection{Pasos Concretos para la Transformación}

\begin{enumerate}[leftmargin=*]
    \item \aplicacion{[Primera acción específica]:} [Describir detalladamente qué debe hacer el creyente]
    \item \aplicacion{[Segunda acción específica]:} [Explicar el segundo paso práctico]
    \item \aplicacion{[Tercera acción específica]:} [Desarrollar la tercera aplicación concreta]
\end{enumerate}

\begin{mdframed}[style=AlertBox]
\textbf{Pregunta para Reflexión Personal:}

\textit{[Formular una pregunta profunda y penetrante que invite a la reflexión personal seria, el autoexamen, y la aplicación individual del mensaje del texto]}
\end{mdframed}

\subsection{Oración de Aplicación}

\begin{center}
\textcolor{forestgreen}{\Large\textbf{Oración de Compromiso}}
\end{center}

\textit{[Incluir aquí una oración extensa y específica que ayude al lector a aplicar los principios del texto a su vida. La oración debe:
- Adorar a Dios por las verdades reveladas en el texto
- Confesar áreas donde no hemos vivido estos principios
- Pedir perdón por fallos específicos
- Solicitar fortaleza divina para vivir estas verdades
- Comprometerse a acciones concretas
- Buscar la transformación del Espíritu Santo]}

\textcolor{olivegreen}{\textbf{En el nombre de Jesucristo, Amén.}}

% ====================
% APÉNDICE
% ====================
\newpage
\appendix

\section{Recursos Utilizados}

\subsection{Herramientas de Estudio Bíblico}

\begin{itemize}[leftmargin=1.5em]
    \item Concordancia Strong (Versión Exhaustiva)
    \item Diccionario Expositivo Vine de Palabras del Antiguo y Nuevo Testamento
    \item Comentario Exegético 1 -- Autor y Editorial
    \item Comentario Exegético 2 -- Autor y Editorial  
    \item Software Bíblico Utilizado -- Versión específica
    \item Léxicos y Gramáticas de Griego y Hebreo
    \item Atlas Bíblico Histórico
\end{itemize}

\subsection{Bibliografía Consultada}

\begin{itemize}[leftmargin=1.5em]
    \item Libro/Artículo 1 -- Referencia Completa
    \item Libro/Artículo 2 -- Referencia Completa
    \item Recurso en Línea -- URL y Fecha de Consulta
    \item Artículo de Revista Teológica -- Referencia Completa
    \item Tesis o Disertación -- Referencia Completa
\end{itemize}

\section{Notas Adicionales}

\subsection{Para Estudio Personal Adicional}

\begin{itemize}[leftmargin=*]
    \item \textbf{Sugerencia de Estudio 1:} [Descripción detallada de cómo profundizar en este tema]
    \item \textbf{Sugerencia de Estudio 2:} [Recomendación para estudio complementario]
    \item \textbf{Recurso Recomendado:} [Libro o material específico para continuar el crecimiento]
    \item \textbf{Pasajes Relacionados:} [Lista de textos para estudio futuro]
\end{itemize}

\subsection{Información de Contacto}

Para preguntas adicionales sobre este estudio o para discusión teológica:

\begin{itemize}[leftmargin=*]
    \item \textbf{Email:} [tu correo electrónico]
    \item \textbf{Teléfono:} [tu número telefónico]
    \item \textbf{Oficina:} [dirección física si aplica]
    \item \textbf{Horarios de Consulta:} [disponibilidad para reuniones]
\end{itemize}

\vspace{2cm}
\begin{center}
\textcolor{forestgreen}{\large\textbf{Ministerio ``La Gloria es del Señor''}}\\
\textcolor{olivegreen}{Escuela de Ministerio Doctrinal}\\
\textcolor{sagegreen}{``Toda la Escritura es inspirada por Dios, y útil para enseñar,\\para redargüir, para corregir, para instruir en justicia'' -- 2 Timoteo 3:16}
\end{center}

\end{document}